ટ્યુરિંગ મશીન તેના સૂચનાગણથી નક્કી થાય છે. સૂચનાગણ પંચકોનો સાન્ત ગણ છે.
દરેક પંચક આપેલી અવસ્થા અને વાંચેલા ચિહ્ન માટે નવી અવસ્થા, લખવાનું ચિહ્ન
અને શીર્ષની ગતિ આપે છે. પંચકોના સાન્ત ગણો ગણનીય હોવાથી દરેક ટ્યુરિંગ
મશીનના સૂચનાગણ સાથે સંખ્યા સાંકળી શકાય છે. એવી કોઈ નિશ્ચિત પદ્ધતિથી તેના
સૂચનાગણને અપાયેલી સંખ્યા તે મશીનનો \emph{સૂચકાંક} છે. આમ સૂચકાંકો વિશે
વાત કરીને આપણે આડકતરી રીતે ટ્યુરિંગ મશીનો વિશે વાત કરી શકીએ છીએ.

ટ્યુરિંગ મશીનના વર્તન અંગેનો એક અગત્યનો પ્રશ્ન એ છે કે તે અંતે અટકશે કે
કેમ. $h(e,n)$ એવું ફલન છે કે જો સૂચકાંક $e$ ધરાવતું ટ્યુરિંગ મશીન
$n$ નિવેશથી શરૂ કરતાં અટકે તો તેનું મૂલ્ય $1$ છે, અને બાકીના કિસ્સામાં
$0$ છે. તેને \emph{અટકવાનું ફલન} કહે છે. એ ફલન પોતે
ટ્યુરિંગ-સંગણનીય છે કે કેમ એ પ્રશ્નને \emph{અટકવાની સમસ્યા} કહે છે.
અટકવાની સમસ્યા ઉકેલી શકાતી નથી. વિકર્ણ દલીલથી આ પરિણામ સ્થાપિત થાય છે.

અટકવાની સમસ્યા એ «ટ્યુરિંગ મશીનોથી $X$ કરી શકાય કે કેમ» એ સ્વરૂપની
સમસ્યાઓના મોટા વર્ગનું એક ઉદાહરણ છે. તર્કમાં બીજી કેન્દ્રીય સમસ્યા
\emph{પ્રથમ-ક્રમ તર્કની નિર્ણય સમસ્યા} છે. આપેલું વાક્ય માન્ય છે કે કેમ
તે નક્કી કરતું ટ્યુરિંગ મશીન છે કે કેમ તે તેનો પ્રશ્ન છે. આ જાણીતી
સમસ્યા ઉકેલી શકાતી નથી. તે ન્યૂનીકરણની દલીલથી સ્થાપિત થાય છે. દરેક ટ્યુરિંગ
મશીન $M$ અને નિવેશ $w$ સાથે પ્રથમ-ક્રમનું વાક્ય $T(M,w) \lif E(M,w)$
જોડી શકાય છે. આ વાક્ય ત્યારે અને માત્ર ત્યારે જ માન્ય છે જ્યારે $w$
નિવેશથી શરૂ કરેલું $M$ અટકે. નિર્ણય સમસ્યા ઉકેલી શકાતી હોત તો તેનો ઉપયોગ
કરીને અટકવાની સમસ્યા પણ ઉકેલી શકાત.
